\section{Limitations}

This draft has five important limitations and pending checks.

First, the current RBTAD result on \liberocorelt{} is a single-seed evaluation. Beyond the Majority reports three-seed averages, so a final SOTA statement requires matching that protocol with paired seeds and comparable rollout budgets.

Second, the same-pipeline APA baseline is not yet complete. Table~\ref{tab:main} includes the APA number reported by Beyond the Majority as context, but it should not be interpreted as a strict local head-to-head comparison until APA is reproduced under our code, checkpoint, preprocessing, and evaluator stack. We are therefore treating the current 40.0\% result as a strong within-pipeline signal against local BC, not as a finished superiority claim over APA.

Third, the 2$\times$2 contribution split is still being completed. RBTAD combines rare-aware behavior cloning and TCAD; the final paper needs matched BC, Rare-BC, TCAD-only, and RBTAD cells under the same seed, batch size, checkpoint, and evaluation protocol before attributing gains to the interaction between the two terms.

Fourth, we do not have physical robot hardware for the Real-World-LT benchmark. We therefore use \liberospatiallt{} as a simulated long-tail generalization test, but this is not a substitute for the original real-world protocol. The current Spatial-LT evidence is useful mainly as a failure-boundary analysis, not as a second main result.

Fifth, negative instruction construction is task-template based and the margin-success causal chain is not yet fully closed. This is simple and low-cost, but it can create false negatives if two tasks share the same early action. The current mild margin and sampling ratio reduce this risk, but do not eliminate it. The final version should report pre/post margin distributions and correlate per-task margin repair with closed-loop success improvements.

The methods improve task-conditioned action boundaries but do not solve all rare relation tasks. In particular, the wine-bottle-to-rack task remains near zero in the current \liberocorelt{} run. On \liberospatiallt{}, RSDF gives a promising seed 7 gain but only a small seed 13 gain, still leaves tasks 5 and 7 at zero, and can reduce task 6 relative to the baseline. Additional Spatial-LT screens after RSDF further narrow the failure mode: full-parameter L2-SP correction, confusion-gated rare behavior cloning, a five-step corrective pulse, and LLM-only relation-localized correction all fail to produce a consistent two-seed gain. This suggests that future Spatial-LT methods should preserve the baseline action distribution directly while applying relation-specific corrective gradients, rather than relying on parameter anchoring, blanket rare weighting, or module freezing alone.